% Generated draft: requires mathematical and source review.
\documentclass[11pt]{article}
\usepackage[margin=25mm]{geometry}
\usepackage{fontspec}
\setmainfont{Latin Modern Roman}
\usepackage{amsmath,amssymb,mathtools,unicode-math}
\setmathfont{Latin Modern Math}
\usepackage{enumitem,needspace,xurl,hyperref}
\setlength{\parindent}{0pt}
\setlength{\parskip}{5pt}
\setlength{\emergencystretch}{3em}
\newcommand{\N}{\mathbb{N}}
\newcommand{\Z}{\mathbb{Z}}
\newcommand{\Q}{\mathbb{Q}}
\newcommand{\R}{\mathbb{R}}
\newcommand{\C}{\mathbb{C}}
\newcommand{\F}{\mathbb{F}}
\newcommand{\A}{\mathbb{A}}
\newcommand{\PP}{\mathbb{P}}
\newcommand{\E}{\mathbb{E}}
\newcommand{\eps}{\varepsilon}
\newcommand{\dd}{\,\mathrm{d}}
\begin{document}
\section*{UnsolvedMath ID: 30006338}
\textbf{Source problem code:} catalogue entry
\subsection{Definitions and mathematical statement}
Write \(\F_p\) for the field with \(p\) elements, where \(p\) is prime, and \(\Z_p^\times\) for the group of units of the ring of \(p\)-adic integers. Put
\[
K=\F_p((u_{n-1})),\qquad K^{\mathrm{perf}}=\varinjlim_{x\mapsto x^p}K.
\]
Here \(K\) is a field of formal Laurent series. Write \(\widehat{K^{\mathrm{perf}}}\) provisionally for the completed perfection mentioned in the source. The completion topology and its compatibility with the group action are not specified in the supplied excerpt. The groups denoted by \(G_m\), their topologies, and their actions on the coefficient objects are also not defined there. For a topological group \(G\) and a specified continuous \(G\)-module \(M\), write \(H^*_{\mathrm{cts}}(G,M)\) for continuous group cohomology; the precise coefficient category and cohomological conventions required here remain unspecified.

\textbf{Supplied comparison and requested extension.} The excerpt reports that, for \(t=n-1\), there is a canonical graded cohomology isomorphism, provisionally transcribed as
\[
H^*_{\mathrm{cts}}\!\left(G_n,\widehat{K^{\mathrm{perf}}}\right)
\cong
H^*_{\mathrm{cts}}\!\left(G_{n-1}\times\Z_p^\times,\F_p\right).
\]
This transcription follows the accompanying prose: the extracted display has damaged superscripts and subscripts, including the completed-perfection notation and the second group factor. It requires comparison with the original typeset formula.

The first requested objective is to remove completed perfection: construct, for every admissible choice of the source parameters, a canonical isomorphism
\[
H^*_{\mathrm{cts}}\!\left(G_n,\F_p((u_{n-1}))\right)
\cong
H^*_{\mathrm{cts}}\!\left(G_{n-1}\times\Z_p^\times,\F_p\right)
\qquad (t=n-1),
\]
with the groups, coefficient actions, topology, and meaning of canonical identification intended in the source. The excerpt does not establish that this isomorphism should be induced by the coefficient map \(K\to\widehat{K^{\mathrm{perf}}}\), so that stronger requirement is not imposed here.

The second requested objective is to extend the mixed-chromatic comparison of the source's Conjecture 2 to every admissible \(t\). Conjecture 2 itself is absent from the supplied record. Consequently, its general coefficient objects, comparison groups, parameter ranges, and comparison maps cannot be stated precisely from this evidence. The excerpt calls the project active work in progress, describes \(t=n-1\) as the easiest case, and expresses an expectation that the methods apply to general \(t\); it does not supply a theorem for general \(t\).

The preceding fragment attributes a comparison to a two-towers construction with torsor maps and a canonical equivalence of quotient \(v\)-stacks. Its diagram and definitions are omitted. It describes a space obtained from projective \((n-1)\)-space by removing all \(\Q_p\)-rational hyperplanes, but does not supply the ambient base field. These contextual assertions do not fill the missing hypotheses above. This is therefore a statement-only draft with explicit unresolved formulation data.
\subsection{Short English statement}
The problem asks for the mixed-chromatic cohomology comparison for every admissible \(t\), and for removal of completed perfection from the reported \(t=n-1\) comparison. The supplied record labels the broader problem partially solved. Missing group definitions, coefficient actions, parameter ranges, and the text of Conjecture 2 prevent a fully self-contained formulation.
\subsection{Sources}
\begin{itemize}
\item Supplied local UnsolvedMath registry record, permanent ID \(30006338\), title ``Mixed-Chromatic Cohomology Beyond the Extreme Case,'' source problem code \texttt{catalogue entry}, in the Oberwolfach Reports Open Problems collection. No individual proposer is identified; the supplied proposed year is 2025.
\item \emph{Non-Archimedean Geometry and Applications} (2025), Oberwolfach Reports, pp.~1517--1578, DOI \texttt{10.4171/owr/2025/29}. Supplied report reference and literature source: \url{https://doi.org/10.4171/owr/2025/29}. No location within the report beyond this page range is supplied.
\item Supplied literature source: \emph{On the rationalization of the K(n)-local sphere}, \url{https://arxiv.org/abs/2402.00960}. The record supplies no authors or precise result location for this reference.
\item Literature assessment dated exactly \texttt{2026-08-22}: the supplied curation labels the problem partially solved, reports that the \(t=n-1\) comparison with completed perfection is proved, and reports that removing completed perfection and proving the general comparison remain unresolved. This is supplied dated curation, not an independently verified literature assessment. The background explicitly requests independent verification of time-sensitive status and cited claims.
\item The registry supplies the statement-status label \texttt{corrected\_verified} and a verification note asserting that the official report was checked and notation repaired. These are supplied provenance claims; no independent source inspection or mathematical verification is asserted in this draft.
\end{itemize}
\end{document}
